\section{Skończone rozkłady i dwie kopie zbioru}
\subsection{Działanie grupy i dozwolone ruchy}
\begin{definicja}[Działanie grupy]
Grupa $G$ działa na zbiorze $Y$, jeśli każdej parze $(g,x)\in G\times Y$ przyporządkowano $gx\in Y$ tak, że $ex=x$ i $(gh)x=g(hx)$. Każde przekształcenie $x\mapsto gx$ jest bijekcją, z odwrotnością $x\mapsto g^{-1}x$. Dla $C\subseteq Y$ piszemy $gC=\{gx:x\in C\}$.
\end{definicja}
Przykłady to działanie grupy na sobie przez mnożenie z lewej, działanie obrotów na sferze oraz działanie ruchów sztywnych na $\R^3$.
\begin{definicja}[Równoważność przez skończony rozkład]\label{def:rownorozklad}
Niech $E,F\subseteq Y$, gdzie $G$ działa na $Y$. Piszemy $E\equivG F$, jeżeli istnieją skończona partycja $E=E_1\sqcup\cdots\sqcup E_k$ oraz $g_1,\ldots,g_k\in G$ takie, że $F=g_1E_1\sqcup\cdots\sqcup g_kE_k$. Słowo \emph{partycja} oznacza, że części są parami rozłączne i pokrywają cały zbiór. Puste części można usuwać. Gdy grupa ruchów jest ustalona, opuszczamy indeks $G$.
\end{definicja}
Istnienie takiego rozkładu jest równoważne istnieniu bijekcji $f:E\to F$, która na każdej ze skończenie wielu części $E_i$ ma postać $f(x)=g_ix$. Rozłączność obrazów daje różnowartościowość, a pokrycie $F$ --- surjektywność.

\textbf{Przykład (dwie części).} Niech grupa $(\Z,+)$ działa na $\Z$ przez przesunięcia. Zbiór $E=\{0,1\}$ jest równoważny zbiorowi $F=\{0,3\}$: część $\{0\}$ pozostawiamy na miejscu, a część $\{1\}$ przesuwamy o $2$. Obrazy $\{0\}$ i $\{3\}$ są rozłączne i pokrywają $F$. Ten przykład pokazuje, że na różnych częściach mogą działać różne elementy grupy.

\begin{lemat}[Składanie rozkładów]\label{lem:skladanie}
Relacja $\equivG$ jest zwrotna, symetryczna i przechodnia. Złożenie rozkładów na odpowiednio $p$ i $q$ części wymaga co najwyżej $pq$ części.
\end{lemat}
\begin{proof}
Zwrotność daje jedna część i ruch $e$. Dla symetrii bierzemy części $g_iE_i$ zbioru docelowego i ruchy $g_i^{-1}$. Przejdźmy do przechodniości. Niech $f:E\to F$ działa jako $g_i$ na $E_i$, a $h:F\to H$ jako $h_j$ na $F_j$. Dzielimy $E$ na zbiory
\begin{equation}\label{eq:3-1}
 E_{ij}'=E_i\cap g_i^{-1}(F_j).
\end{equation}
Dla każdego $x\in E$ istnieje dokładnie jedno $i$ z $x\in E_i$, a następnie dokładnie jedno $j$ z $g_ix\in F_j$. Zatem zbiory $E_{ij}'$ tworzą partycję. Na $E_{ij}'$ złożenie $h\circ f$ jest pojedynczym ruchem $h_jg_i$. Ponieważ złożenie bijekcji jest bijekcją, obrazy tych części są rozłączne i pokrywają $H$. Liczba par $(i,j)$ wynosi $pq$.
\end{proof}

\subsection{Dwie kopie oznaczone etykietami}
Symbol $E\sqcup E$ oznacza sumę rozłączną dwóch kopii $E$. Będziemy ją przedstawiać jako $2E=E\times\{1,2\}$. Pary $(x,1)$ i $(x,2)$ oznaczają dwa różne miejsca docelowe. Etykieta nie jest dodatkową współrzędną euklidesową.

W rozkładzie między kopiami oznaczonymi etykietami każda część ma jedną ustaloną etykietę źródłową i jedną docelową; przekształcenie ma postać $(x,i)\mapsto(gx,j)$. Etykiety służą jedynie do rozróżnienia kopii. Na końcu umieścimy kopie w różnych miejscach $\R^3$ za pomocą przesunięć.

Dowód lematu~\ref{lem:skladanie} działa bez zmian, gdy w przecięciach \eqref{eq:3-1} uwzględniamy zgodność etykiety pośredniej. Jeżeli $E\equivG F$, możemy wykonać ten sam rozkład osobno w obu kopiach, otrzymując $2E\equivG2F$.
\begin{definicja}[Zbiór paradoksalny]\label{def:paradoksalny}
Zbiór $E$ jest \emph{paradoksalny} dla danego działania grupy, jeśli $E\equivG2E$. Równoważnie, istnieje partycja $E=A_1\sqcup\cdots\sqcup A_m\sqcup B_1\sqcup\cdots\sqcup B_n$ i elementy $g_i,h_j\in G$ takie, że
\[
 E=\bigsqcup_{i=1}^{m}g_iA_i,\qquad E=\bigsqcup_{j=1}^{n}h_jB_j.
\]
Pierwsze równanie odtwarza kopię z etykietą $1$, drugie kopię z etykietą $2$. Źródłowe części nie są kopiowane.
\end{definicja}
Niepusty zbiór paradoksalny musi być nieskończony: dla niepustego zbioru skończonego bijekcja $E\to2E$ wymagałaby $|E|=2|E|$. W dalszej części zbudujemy taki rozkład dla $F_2$, a następnie przeniesiemy go na sferę i kulę.
\begin{wniosek}[Przenoszenie paradoksalności]\label{wn:przenoszenie}
Jeśli $E\equivG F$ i $F\equivG2F$, to $E\equivG2E$.
\end{wniosek}
\begin{proof}
Składamy trzy skończone rozkłady: $E\equivG F\equivG2F\equivG2E$. Ostatnia równoważność jest dwukrotnym zastosowaniem odwrotnego rozkładu $F\equivG E$. Lemat~\ref{lem:skladanie} zapewnia, że po złożeniu nadal mamy skończenie wiele części i pojedynczy dozwolony ruch na każdej.
\end{proof}
\subsection{Przesunięcie wzdłuż nieskończonego łańcucha}
\begin{lemat}[Przesunięcie łańcucha]\label{lem:lancuch}
Niech $D\subseteq E\subseteq Y$, a $T:Y\to Y$ będzie przekształceniem wyznaczonym przez pewien element $G$. Załóżmy, że zbiory $D,T(D),T^2(D),\ldots$ są parami rozłączne i wszystkie zawierają się w $E$. Wówczas $E\equivG E\setminus D$ przy użyciu co najwyżej dwóch części.
\end{lemat}
\begin{proof}
Niech $L=\bigcup_{n=0}^{\infty}T^n(D)$. Rozłączność poziomów daje $T(L)=\bigcup_{n=1}^{\infty}T^n(D)=L\setminus D$. Dzielimy $E$ na dwie części $L$ i $E\setminus L$. Pierwszą przenosimy przez $T$, drugą pozostawiamy bez zmian. Obrazy to odpowiednio $L\setminus D$ i $E\setminus L$; są rozłączne i ich suma wynosi $E\setminus D$.
\end{proof}
\begin{uwaga}
Użyliśmy nieskończonej unii do zdefiniowania jednego kawałka $L$. Nie wykonujemy nieskończenie wielu ruchów na różnych poziomach: na całym $L$ działa jeden i ten sam ruch $T$. Rozkład ma więc dwie części, mimo że jedna z nich jest nieskończoną unią.
\end{uwaga}

